\documentclass{article}
\usepackage[T1]{fontenc}
\usepackage{lmodern}
\usepackage{graphicx}
\usepackage{amsmath,amssymb}
\usepackage[a4paper,margin=2cm]{geometry}
\usepackage[absolute,overlay]{textpos}
\setlength{\TPHorizModule}{1mm}
\setlength{\TPVertModule}{1mm}
\usepackage[indonesian]{babel}
\usepackage{hyperref}
\setlength{\emergencystretch}{2em}
% unit: Halaman 62

\begin{document}

% === Halaman 62 (python/native) ===

62

Varian Vigenere Cipher

Untuk mengatasi serangan dengan metode Kasiski, maka dibuat varian 

Vigenere Cipher sebagai berikut:

1.

Full Vigènere cipher

•

Setiap baris di dalam tabel tidak menyatakan pergeseran huruf, tetapi 

merupakan permutasi huruf-huruf alfabet (lihat contoh table pada 

halaman berikut). 

•

Tabel tersebut harus dirahasiakan.

•

Proses enkripsi dan dekripsi tetap sama seperti Vigenere standard:

\end{document}
